\documentclass[12pt]{article}
\usepackage{graphicx} % Required for inserting images
\usepackage[margin=0.5in]{geometry}
\usepackage{amsmath, amssymb}
\usepackage{mathrsfs}


\usepackage{soul}


% Dr. David Brown Commands
\newcommand{\ds}{\displaystyle}
\newcommand{\R}{\mathbb{R}}
\newcommand{\N}{\mathbb{N}}
\newcommand{\Q}{\mathbb{Q}}
\newcommand{\C}{\mathbb{C}}


% Commands to get script r
\usepackage{calligra}
\DeclareMathAlphabet{\mathcalligra}{T1}{calligra}{m}{n}
\DeclareFontShape{T1}{calligra}{m}{n}{<->s*[2.2]callig15}{}
\newcommand{\scripty}[1]{\ensuremath{\mathcalligra{#1}}}

% Unit commands
\newcommand{\degree}{\text{°}}
\newcommand{\unitSpeed}{\frac{\text{m}}{\text{s}}}
\newcommand{\unitAcceleration}{\frac{\text{m}}{\text{s}^2}}

% Constant commands
\newcommand{\avogadro}{6.022\times10^{23}}

\newcommand{\boltzmannConstK}{1.381\times10^{-23}\frac{\text{J}}{\text{K}}}
\newcommand{\boltzmannConstInEV}{8.617\times10^{-5}\frac{\text{eV}}{\text{K}}}
\newcommand{\planckConst}{6.626\times10^{-34}\text{J}\cdot\text{s}}

\newcommand{\atmP}{1.01325\times10^5\,\text{Pa}}

% Known Value Commands
\newcommand{\electronMass}{9.109\times10^{-31}\,\text{kg}}
\newcommand{\protonMass}{1.673\times10^{-27}\,\text{kg}}

% Commands to easily write partial derivatives
\newcommand{\partDeriv}[2]{\frac{\partial #1}{\partial #2}}
\newcommand{\doublePartDeriv}[2]{\frac{\partial^2 #1}{\partial^2 #2}}


\title{Problem 7.29}
\author{Gabriel Decker}


\begin{document}


\maketitle
\begin{flushleft}



In this problem, we will solve the book's equation for the total energy of a system of fermions with a density of energy states $g(\epsilon)$ (book equation 7.54) for a Fermi gas with $kT\ll\epsilon_F$ in three dimensions.
We will use the Sommerfield Expansion to do so.
We will first go from this equation:
\begin{equation}
    \label{eq:start}
    U=\int_0^\infty\epsilon g(\epsilon)\frac{1}{e^{(\epsilon-\mu)/kT}+1}\,d\epsilon
\end{equation}
to this equation.
\begin{equation}
    U=\frac{3}{5}N\frac{\mu^{5/2}}{\epsilon_F^{3/2}}+\frac{3\pi^2}{8}N\frac{(kT)^2}{\epsilon_F}+\dots
\end{equation}\\~\\

In this case, we get the following equation for $g(\epsilon)$ from the book.
\begin{equation}
    g(\epsilon)=\frac{3N}{2\epsilon_F^{3/2}}\epsilon^{1/2}
\end{equation}
and this can be expressed as $g(\epsilon)=g_0\epsilon^{1/2}$.
Also, recall that the fraction with the exponent in equation \ref{eq:start} is the fermi-dirac distribution.
\begin{equation}
    \overline{n}_\text{FD}(\epsilon)=\frac{1}{e^{(\epsilon-\mu)/kT}+1}
\end{equation}\\~\\

With all that preable, we begin applying the Sommerfield Expansion by plugging in our $g$ and doing integration by parts.
$$U=\int_0^\infty\epsilon g(\epsilon)\overline{n}_\text{FD}(\epsilon)\,d\epsilon$$
$$\implies U=\int_0^\infty\epsilon g_0\epsilon^{1/2}\overline{n}_\text{FD}(\epsilon)\,d\epsilon$$
$$\implies U=\int_0^\infty\epsilon^{3/2}\overline{n}_\text{FD}(\epsilon)\,d\epsilon$$
$dv=g_0\epsilon^{3/2}\,d\epsilon$, so $v=g_0\frac{2}{5}\epsilon^{5/2}$.\\
$u=\overline{n}_\text{FD}(\epsilon)$, so $du=\frac{d\overline{n}_\text{FD}(\epsilon)}{d\epsilon}$.\\
By integration by parts, this implies the following.
$$\implies U=g_0\frac{2}{5}\epsilon^{5/2}\overline{n}_\text{FD}(\epsilon)\Big|_0^\infty-g_0\int_0^\infty\frac{2}{5}\epsilon^{5/2}\frac{d\overline{n}_\text{FD}(\epsilon)}{d\epsilon}\,d\epsilon$$
$$\implies U=g_0\frac{2}{5}\epsilon^{5/2}\overline{n}_\text{FD}(\epsilon)\Big|_0^\infty+\frac{2}{5}g_0\int_0^\infty\epsilon^{5/2}\left(-\frac{d\overline{n}_\text{FD}(\epsilon)}{d\epsilon}\right)\,d\epsilon$$
The first term goes to zero because at $\epsilon\rightarrow\infty$, $\overline{n}_\text{FD}\rightarrow0$ exponentially, while at $\epsilon=0$, $\overline{n}_\text{FD}$ becomes just some number while $\epsilon^{5/2}=0$.
So, we get the following integral left.
$$\implies U=\frac{2}{5}g_0\int_0^\infty\epsilon^{5/2}\left(-\frac{d\overline{n}_\text{FD}(\epsilon)}{d\epsilon}\right)\,d\epsilon$$\\~\\

We know that far from $\epsilon=\mu$, the derivative of $\overline{n}_\text{FD}$ is nearly zero since there is nearly no change in the function, except near $\mu$, especially at lower temperatures, which we're assuming.
Therefore, our approximation is still pretty accurate if we extend the bounds of integration to negative infinity.
This makes our integrals more symmetric.
$$\implies U=\frac{2}{5}g_0\int_{-\infty}^\infty\epsilon^{5/2}\left(-\frac{d\overline{n}_\text{FD}(\epsilon)}{d\epsilon}\right)\,d\epsilon$$\\~\\

The book provided the following solution to the derivative in the integrand.
$$-\frac{d\overline{n}_\text{FD}(\epsilon)}{d\epsilon}=\frac{1}{kT}\frac{e^x}{(e^x+1)^2}$$
with $x=(\epsilon-\mu)/kT$ and thus $dx=d\epsilon/kT$.
We can put this in the integral and get the following.
$$\implies U=\frac{2}{5}g_0\int_{-\infty}^\infty\epsilon^{5/2}\left(\frac{1}{kT}\frac{e^x}{(e^x+1)^2}\right)\,d\epsilon$$
$$\implies U=\frac{2}{5}g_0\int_{-\infty}^\infty\frac{e^x}{(e^x+1)^2}\epsilon^{5/2}\frac{d\epsilon}{kT}$$
$$\implies U=\frac{2}{5}g_0\int_{-\infty}^\infty\frac{e^x}{(e^x+1)^2}\epsilon^{5/2}\,dx$$
We still need to express $\epsilon^{5/2]}$ in terms of $x$, but this can be simplified by taking the taylor expansion of it, centered around $\mu$.
This expansion is the following.
$$\epsilon^{5/2}=\mu^{5/2}+\frac{5}{2}\mu^{3/2}(\epsilon-\mu)+\frac{15}{8}\mu^{1/2}(\epsilon-\mu)^2+\dots$$
Since $(\epsilon-\mu)=xkT$, we can substitute that in here.
$$\implies\epsilon^{5/2}=\mu^{5/2}+\frac{5}{2}\mu^{3/2}xkT+\frac{15}{8}\mu^{1/2}(xkT)^2+\dots$$\\~\\

We can now plug this into our integral to get only $x$ dependence.
$$\implies U=\frac{2}{5}g_0\int_{-\infty}^\infty\frac{e^x}{(e^x+1)^2}\left[\mu^{5/2}+\frac{5}{2}\mu^{3/2}xkT+\frac{15}{8}\mu^{1/2}(xkT)^2+\dots\right]\,dx$$
This integral can be solved using methods similar to the book, but let's try doing it using sympy.





\end{flushleft}

\end{document}
